\documentclass[a4paper,12pt]{article}
\usepackage[utf8]{inputenc}
\usepackage[english]{babel}
\usepackage{amsmath, amssymb}
\usepackage{enumitem}
\usepackage{geometry}
\geometry{margin=2.5cm}

\title{Advanced Quantum Mechanics - Comprehensive Quiz}
\author{Based on TU Darmstadt Lecture Material 2025/2026}
\date{}

\begin{document}

\maketitle

\section*{Instructions}
This quiz covers: Fundamentals (F), Scattering Theory (S), Many Body Physics (M), and Relativistic Quantum Mechanics (R). Each section contains 5 blocks with 5 statements each. Statements are either \textbf{True} or \textbf{False}. The number of correct statements is indicated for each block.

\newpage

\section{Fundamentals of Quantum Mechanics (F)}

\subsection*{Block F1: Schrödinger Equation and Dirac Notation}
\textit{Number of correct statements: 4}
\begin{enumerate}[label=\roman*.]
    \item The state vector $|\psi(t)\rangle$ lives in a complex Hilbert space and contains all physical information.
    \item The inner product $|\langle\phi|\psi\rangle|^2$ represents the probability of finding a system in state $|\phi\rangle$.
    \item For an operator to represent a measurable observable, it must be Unitary.
    \item The TDSE ensures that the norm of the state is preserved if the Hamiltonian is Hermitian.
    \item In Dirac notation, a bra vector $\langle\psi|$ corresponds to the Hermitian conjugate of the ket $|\psi\rangle$.
\end{enumerate}

\subsection*{Block F2: Time Evolution and Dyson Series}
\textit{Number of correct statements: 3}
\begin{enumerate}[label=\roman*.]
    \item For a time-independent Hamiltonian, the evolution operator is $\hat{U}(t) = e^{-i\hat{H}t/\hbar}$.
    \item If $[\hat{H}(t_1), \hat{H}(t_2)] \neq 0$, the solution is a simple exponential of the time-integrated Hamiltonian.
    \item The time-ordering operator $\mathcal{T}$ ensures earlier times act to the left of later times.
    \item The Dyson expansion is a power series representation of the time-ordered exponential.
    \item In the energy eigenbasis, time evolution attaches a phase factor $e^{-iE_n t/\hbar}$ to each component.
\end{enumerate}

\subsection*{Block F3: Two-Level Systems (Qubits)}
\textit{Number of correct statements: 4}
\begin{enumerate}[label=\roman*.]
    \item A Hamiltonian $H = \hbar \Omega \sigma_x$ can cause transitions between $|0\rangle$ and $|1\rangle$.
    \item The eigenvalues of $\sigma_x$ are $\pm 1$, leading to energy levels $\pm \hbar \Omega$.
    \item The state $|+\rangle = \frac{1}{\sqrt{2}}(|0\rangle + |1\rangle)$ is an eigenvector of $\sigma_z$.
    \item Time evolution of $|0\rangle$ under $H = \hbar \Omega \sigma_x$ gives $|\psi(t)\rangle = \cos(\Omega t)|0\rangle - i \sin(\Omega t)|1\rangle$.
    \item $\sigma_x$ is its own inverse and is also Hermitian.
\end{enumerate}

\subsection*{Block F4: Position and Momentum Representations}
\textit{Number of correct statements: 2}
\begin{enumerate}[label=\roman*.]
    \item The momentum operator in position representation is $\hat{p} \to -i\hbar \nabla$.
    \item The completeness relation $\int dx |x\rangle\langle x| = 1$ allows defining $\psi(x) = \langle x | \psi \rangle$.
    \item The position operator $\hat{x}$ acts as a derivative in the position representation.
    \item The commutator $[\hat{x}, \hat{p}]$ is equal to zero in the position representation.
    \item Momentum eigenstates are localized Delta functions in position space.
\end{enumerate}

\subsection*{Block F5: The Free Particle}
\textit{Number of correct statements: 3}
\begin{enumerate}[label=\roman*.]
    \item For a free particle ($V(x)=0$), the energy spectrum is continuous and positive ($E > 0$).
    \item Plane wave solutions $\phi_k(x) = \frac{1}{\sqrt{2\pi}}e^{ikx}$ are square-integrable ($L^2$).
    \item The dispersion relation for a non-relativistic free particle is $E_k = \frac{\hbar^2 k^2}{2m}$.
    \item Any initial state $\psi(x,0)$ can be evolved by applying time-dependent phases to its Fourier components.
    \item Negative energy solutions for a free particle are standard bound states in non-relativistic QM.
\end{enumerate}

\newpage

\section{Scattering Theory (S)}

\subsection*{Block S1: Basics and Cross-Sections}
\textit{Number of correct statements: 4}
\begin{enumerate}[label=\roman*.]
    \item The scattering amplitude $f(\theta, \phi)$ relates theory to experimental cross-sections.
    \item The differential cross-section is defined as $d\sigma/d\Omega = |f(\theta, \phi)|^2$.
    \item Scattering is used to probe the shape and strength of interaction potentials.
    \item Integrating the differential cross-section $|f|^2$ over the full solid angle yields the total cross-section.
    \item The scattering amplitude $f$ has the physical dimensions of energy.
\end{enumerate}

\subsection*{Block S2: Lippmann-Schwinger Equation}
\textit{Number of correct statements: 3}
\begin{enumerate}[label=\roman*.]
    \item The Lippmann-Schwinger equation is an integral form of the Schrödinger equation.
    \item It describes $|\psi^{(+)}\rangle$ as a sum of the incident plane wave and a scattered term.
    \item The Green's function $\hat{G}(E_i)$ uses an $i\epsilon$ term to ensure outgoing wave boundary conditions.
    \item The Lippmann-Schwinger equation is only valid for potentials stronger than the kinetic energy.
    \item In the momentum basis, the free Green's function operator is non-diagonal.
\end{enumerate}

\subsection*{Block S3: Born Approximation}
\textit{Number of correct statements: 3}
\begin{enumerate}[label=\roman*.]
    \item The first Born approximation replaces the full state with the incident plane wave in the integral.
    \item The Born approximation is most accurate for very strong potentials at low energies.
    \item For a spherically symmetric potential, the Born amplitude depends only on the momentum transfer $q$.
    \item The Born amplitude is proportional to the Fourier transform of the potential $V(r)$.
    \item The Born approximation always strictly preserves the unitarity of the S-matrix.
\end{enumerate}

\subsection*{Block S4: Partial Wave Expansion}
\textit{Number of correct statements: 4}
\begin{enumerate}[label=\roman*.]
    \item Partial wave analysis decomposes scattering into components of definite angular momentum $l$.
    \item At low energies ($ka \ll 1$), the $s$-wave ($l=0$) usually dominates the scattering.
    \item The phase shift $\delta_l$ characterizes the shift of the radial wavefunction's nodes due to $V(r)$.
    \item If a phase shift $\delta_l$ is zero, there is no scattering in that specific $l$-channel.
    \item For a hard sphere, all phase shifts $\delta_l$ are always positive.
\end{enumerate}

\subsection*{Block S5: Low Energy Limits (Hard Sphere)}
\textit{Number of correct statements: 2}
\begin{enumerate}[label=\roman*.]
    \item In the low-energy limit, scattering from short-range potentials becomes isotropic.
    \item The scattering length $a$ is defined by the $k \to 0$ limit of the $s$-wave amplitude.
    \item The Born approximation is always exact for a hard sphere potential at any energy.
    \item For a square-well potential, scattering is always anisotropic as $k \to 0$.
    \item The total cross-section $\sigma$ at $k \to 0$ for a hard sphere of radius $R$ is exactly $\pi R^2$.
\end{enumerate}

\newpage

\section{Many Body Physics (M)}

\subsection*{Block M1: Indistinguishable Particles}
\textit{Number of correct statements: 3}
\begin{enumerate}[label=\roman*.]
    \item Distinguishable particles require wavefunctions to be totally symmetric.
    \item Bosons have symmetric wavefunctions and follow Bose-Einstein statistics.
    \item Fermions follow the Pauli Exclusion Principle (antisymmetric wavefunctions).
    \item Exchange symmetry of a many-body wavefunction has no effect on physical observables.
    \item Indistinguishability is a fundamental property of identical particles like electrons.
\end{enumerate}

\subsection*{Block M2: Second Quantization - Operators}
\textit{Number of correct statements: 4}
\begin{enumerate}[label=\roman*.]
    \item Bosonic creation/annihilation operators satisfy $[\hat{a}_i, \hat{a}_j^\dagger] = \delta_{ij}$.
    \item Fermionic operators satisfy anti-commutation relations $\{\hat{c}_i, \hat{c}_j^\dagger\} = \delta_{ij}$.
    \item For fermions, $(\hat{c}_j^\dagger)^2 = 0$, representing the Pauli principle.
    \item The number operator $\hat{n}_j = \hat{c}_j^\dagger \hat{c}_j$ counts particles in state $j$.
    \item Applying an annihilation operator to the vacuum state $|0\rangle$ yields $|0\rangle$.
\end{enumerate}

\subsection*{Block M3: Fock States and Distributions}
\textit{Number of correct statements: 3}
\begin{enumerate}[label=\roman*.]
    \item A Fock state $|n_0, n_1, ...\rangle$ is an eigenstate of the mode number operators.
    \item In second quantization, the vacuum state $|0\rangle$ is the state with zero particles.
    \item A thermal state is a coherent superposition with a single fixed phase.
    \item For bosons, any number of particles can occupy the same mode $j$.
    \item The mean occupation number in a thermal state follows the Bose-Einstein distribution for bosons.
\end{enumerate}

\subsection*{Block M4: Gross-Pitaevskii and Hubbard Models}
\textit{Number of correct statements: 4}
\begin{enumerate}[label=\roman*.]
    \item The Gross-Pitaevskii equation (GPE) is a non-linear Schrödinger equation for BECs.
    \item The GPE is derived by taking the expectation value of the Heisenberg field equations.
    \item The Hubbard model describes hopping ($t$) and on-site interactions ($U$) on a lattice.
    \item In the limit $U \gg t$ at half-filling, the Hubbard model maps to the Heisenberg spin model.
    \item The Hubbard model is only applicable to non-interacting bosonic systems.
\end{enumerate}

\subsection*{Block M5: Magnetism and Spin Models}
\textit{Number of correct statements: 2}
\begin{enumerate}[label=\roman*.]
    \item The Heisenberg model describes interactions between neighboring spins $\mathbf{S}_i \cdot \mathbf{S}_j$.
    \item The Ising model is a case where spins align only along a single axis (e.g., $z$).
    \item In a transverse field Ising model, the system is always ferromagnetic for any field strength.
    \item The Jordan-Wigner transformation maps 1D spin systems to non-interacting bosons.
    \item At the critical point $h=J$ in the transverse Ising model, the energy gap remains large.
\end{enumerate}

\newpage

\section{Relativistic Quantum Mechanics (R)}

\subsection*{Block R1: Lorentz Invariance}
\textit{Number of correct statements: 4}
\begin{enumerate}[label=\roman*.]
    \item Non-relativistic QM fails when kinetic energy is comparable to rest mass energy.
    \item Physical laws must be invariant under Lorentz transformations in relativistic theory.
    \item The spacetime interval $ds^2$ is an invariant scalar in Minkowski space.
    \item The Schrödinger equation is relativistically covariant because it is first-order in time.
    \item Proper time $\tau$ is a Lorentz scalar representing time in the particle's rest frame.
\end{enumerate}

\subsection*{Block R2: Metric and Four-Vectors}
\textit{Number of correct statements: 4}
\begin{enumerate}[label=\roman*.]
    \item The metric tensor $g_{\mu\nu}$ is used to raise and lower four-vector indices.
    \item Lorentz transformations must satisfy $\Lambda^T g \Lambda = g$.
    \item Common Minkowski signatures are $(1, -1, -1, -1)$ or $(-1, 1, 1, 1)$.
    \item The Minkowski metric in Cartesian coordinates must have non-zero off-diagonal terms.
    \item The four-momentum $p^\mu$ is $(E/c, \mathbf{p})$.
\end{enumerate}

\subsection*{Block R3: Dirac Equation}
\textit{Number of correct statements: 3}
\begin{enumerate}[label=\roman*.]
    \item The Dirac equation is first-order in both time and space derivatives.
    \item It was designed to ensure a positive-definite probability density.
    \item The Dirac equation naturally predicts spin-1/2 for the described particles.
    \item The $\alpha$ and $\beta$ matrices in the Dirac Hamiltonian are $2 \times 2$ matrices.
    \item The Dirac equation only applies to massless particles.
\end{enumerate}

\subsection*{Block R4: Spinors and Solutions}
\textit{Number of correct statements: 4}
\begin{enumerate}[label=\roman*.]
    \item Dirac wavefunctions are four-component spinors.
    \item For $\mathbf{p}=0$, there are two $+mc^2$ energy solutions and two $-mc^2$ solutions.
    \item In the non-relativistic limit, two spinor components are "large" and two are "small".
    \item Negative energy solutions can be ignored as they have no physical meaning in Dirac theory.
    \item The Feynman-Stueckelberg interpretation treats negative energy states as antiparticles.
\end{enumerate}

\subsection*{Block R5: Klein Paradox}
\textit{Number of correct statements: 4}
\begin{enumerate}[label=\roman*.]
    \item The Klein paradox involves high transmission through a barrier as the potential height increases.
    \item It occurs when the potential step $V$ exceeds $E + mc^2$.
    \item The paradox is resolved by considering particle-antiparticle pair production at the step.
    \item In the Klein zone, the wavefunction inside the barrier becomes oscillatory.
    \item Relativistic tunneling is identical to the exponential decay seen in Schrödinger tunneling.
\end{enumerate}

\newpage

\section*{Solutions}

\subsection*{Fundamentals (F)}
\begin{itemize}
    \item \textbf{F1:} i, ii, iv, v (4)
    \item \textbf{F2:} i, iv, v (3)
    \item \textbf{F3:} i, ii, iv, v (4)
    \item \textbf{F4:} i, ii (2)
    \item \textbf{F5:} i, iii, iv (3)
\end{itemize}

\subsection*{Scattering (S)}
\begin{itemize}
    \item \textbf{S1:} i, ii, iii, iv (4) --- \textit{iv is True: $\int |f|^2 d\Omega = \sigma_{tot}$.}
    \item \textbf{S2:} i, ii, iii (3)
    \item \textbf{S3:} i, iii, iv (3)
    \item \textbf{S4:} i, ii, iii, iv (4)
    \item \textbf{S5:} i, ii (2) --- \textit{v is False: $\sigma_{k \to 0} = 4\pi R^2$.}
\end{itemize}

\subsection*{Many Body (M)}
\begin{itemize}
    \item \textbf{M1:} ii, iii, v (3)
    \item \textbf{M2:} i, ii, iii, iv (4)
    \item \textbf{M3:} i, iv, v (3)
    \item \textbf{M4:} i, ii, iii, iv (4) --- \textit{iv is True: Hubbard $\to$ Heisenberg at $U \gg t$.}
    \item \textbf{M5:} i, ii (2)
\end{itemize}

\subsection*{Relativistic (R)}
\begin{itemize}
    \item \textbf{R1:} i, ii, iii, v (4) --- \textit{iv is False: Schrödinger is not covariant.}
    \item \textbf{R2:} i, ii, iii, v (4) --- \textit{iv is False: $g_{\mu\nu}$ is diagonal.}
\end{itemize}

\begin{itemize}
    \item \textbf{R3:} i, ii, iii (3) --- \textit{iv is False ($4 \times 4$ matrices), v is False (Massive particles).}
    \item \textbf{R4:} i, ii, iii, v (4) --- \textit{iv is False: Negative energy is essential.}
    \item \textbf{R5:} i, ii, iii, iv (4) --- \textit{v is False: Different from non-relativistic tunneling.}
\end{itemize}

\end{document}